\section{Discussion and limitations}
The study separates architecture, calibration, and allocation. An identity-anchored correction provides a fair neural comparison with persistence, but its target has overlapping observations. The window-update forecast explains much of the gain without a network. Claims about latent or forward variance require evidence beyond low error on a smoothed rolling window.

Fitting Heston parameters does not remove model risk. Predictive likelihood gains are mixed, despite tighter numerical searches and convergence of all refinement stages. The estimates remain conditional on a daily Euler likelihood, a Feller restriction, finite particles, and a bounded search. Resampling can produce irregular likelihood surfaces and weakly identified parameters. Displayed estimates have only two significant figures, and numerical likelihood variation is not parameter-estimation uncertainty.

With four risky assets and a quadratic risk term, allocations can change composition at fixed gross exposure. The cash-enabled comparison demonstrates that cash preference may nevertheless remain under estimated premia and the selected risk coefficient. A diversified decision space is necessary to study relative allocation, but it does not ensure attractive investment opportunities.

The 20,000-scenario experiment makes the numerical comparison more informative. Independent-seed allocation variation is smaller than the average robustness tilt, and the direction of the 2025 adaptive-minus-nominal return difference is stable across the three seeds. Statistical intervals nevertheless include zero. Numerical resolution and evidence of economic value are distinct questions. An absence of a statistically reliable gain under these settings is a substantive result, rather than a claim that every robustness effect is hidden by the simulator.

The 2025 confirmation sample is unusually favorable to gold. GLD's adjusted-close calendar return is 63.7\%, compared with 17.7\% for SPY, 12.7\% for IWM, and 4.2\% for TLT. The historical mean--CVaR benchmark averages about 90\% GLD weight and an effective asset count of 1.21. Its roughly 61\% annualized CAGR is largely a concentrated gold position. It is not broad evidence that its risk model forecasts well. Calendar close-to-close asset returns and the 247-day adjusted-open portfolio slice have different boundaries, so these figures are explanatory context rather than an exact return attribution.

Other limitations remain material. SPY and IWM share substantial equity exposure, and four ETFs cannot represent all market opportunities. Dependence uses shrunk return correlation and a restricted shock construction. Neural rolling-volatility forecasts approximate an instantaneous variance state. Drifts are noisy and do not establish return predictability. Adjusted-open execution assumes immediate fills after observing the mark; spreads, slippage, taxes, and delay are omitted. The linear fee proxy differs from exact compounded deductions by interaction terms.

The exploratory sample motivated the design, and confirmation covers only one historical year. Bootstrap intervals assume a useful degree of stationarity, do not refit models, and are unadjusted for multiple comparisons. Exact replication also requires access to matching third-party inputs. Public code, hashes, model records, and summaries make the computational procedure inspectable, but they do not grant data rights or ensure that providers retain an unchanged historical snapshot.

Future work can use intraday measures or direct forward targets, assess parameter identification with more accurate likelihoods, test other universes and confirmation periods, and calibrate transport radii against held-out distributional errors.

\section{Conclusion}
This selective review combines established DRO methods with residual volatility forecasting, fitted physical Heston dynamics, and convex multi-asset allocation. Persistence is explicit in the neural architecture, structural parameters use historically available returns, and fixed total ETF exposure separates composition from cash-based de-risking.

Residual forecasting improves on persistence while a simple window-update model remains competitive. Heston likelihood gains are mixed. Higher-scenario paired allocation runs resolve a small robustness-induced change in weights, but do not establish a statistically reliable return advantage over matched nominal allocation. Strong 2025 CAGRs must be interpreted alongside a gold rally and concentrated gold exposure. The public computational record supports a reproducible case study of these established components, with conditional data access and restrained performance claims.

\section*{Code and data availability}
Code, configuration, 168 trained network checkpoints, 28 annual parameter-fit records, aggregate results, input hashes, and figure/table scripts are available in the versioned public GitHub release \cite{repository}:\par
\noindent{\footnotesize\url{https://github.com/Sheehan007/Distributionally-Robust-Optimization/releases/tag/v1.0.0}}\par
Provider observations, raw API responses, and date-by-date market and portfolio caches are excluded from public artifacts. The release documents Yahoo's and Cboe's terms and provides a downloader requiring explicit acknowledgment of those terms. Exact numerical replication is conditional on appropriate access and matching canonical input hashes. A provider-history change produces a new-data evaluation, not an exact replication of the reported snapshot.

\section*{Author and tool disclosure}
Sheehan Mathur is the sole author. Anthropic Claude assisted with Table~\ref{tab:families} and parts of the literature-review text. OpenAI Codex assisted with implementation audits, experiment development and execution, drafting, reference checking, and typesetting. The author is responsible for the final scientific interpretation and submission.
